% GENERATED FILE. Edit canonical lessons, then run npm run generate:books.
% canonical-source-hash: fnv1a64:17a81a89f1f8a92f
% canonical-lessons: PT-C289-fraude, PT-C289-burla, PT-C289-burlar, PT-C289-burlao, PT-C289-cumplice

\chapter{Burlas e fraudes}
\label{ch:pt-burlas-e-fraudes}
\hlchaptermodality{289}

\begin{chapteropening}
\textbf{By the end of this chapter:} \emph{I can talk about fraud and scams, the con man and his accomplice.}

Complete the last lesson of chapter 289: I can talk about fraud and scams, the con man and his accomplice.
\end{chapteropening}

\section[a fraude]{a fraude — the fraud}

\label{lesson:PT-C289-fraude}



\begin{quote}
Before the new one: say the Portuguese for to break open, then the Portuguese for the break-in.
\end{quote}

\subsection*{You'll want to know: a fraude}
\textbf{a fraude} — \textquotedblleft{}the fraud\textquotedblright{}.

\begin{quote}
A polícia descobriu uma fraude com cartões bancários.
\end{quote}

\begin{quote}
\emph{The police uncovered a fraud with bank cards.}
\end{quote}

\emph{Fraude fiscal} — tax fraud.

\begin{cousinweb}
\emph{fraude} is a learned borrowing from Latin \emph{fraus, fraudis} \textquotedblleft{}deceit, harm\textquotedblright{} (Lewis \& Short).
\end{cousinweb}

\begin{grammarlens}[title={What you've built}]
One more word for telling, explaining and describing at B1.
\end{grammarlens}

\subsection*{Your turn}
\begin{itemize}
\raggedright
  \item \emph{Say it:} \emph{a fraude}
  \item \emph{Say it:} \emph{a fraude}, once more
  \item \emph{Say it:} a sentence of your own with \emph{a fraude}
  \item \emph{Recall:} say the Portuguese for to break open, then the Portuguese for the break-in, then say \emph{a fraude} again
\end{itemize}

\begin{tcolorbox}[breakable,colback=teal!4,colframe=teal!35!black,title={Before you move on}]
What does a fraude mean? (\textquotedblleft{}The fraud\textquotedblright{}.) Say it once more.
\end{tcolorbox}
\section[a burla]{a burla — the scam, the swindle}

\label{lesson:PT-C289-burla}



\begin{quote}
Before the new one: say the Portuguese for the break-in, then the Portuguese for the fraud.
\end{quote}

\subsection*{You'll want to know: a burla}
\textbf{a burla} — \textquotedblleft{}the scam, the swindle\textquotedblright{}.

\begin{quote}
Aquele e-mail era uma burla.
\end{quote}

\begin{quote}
\emph{That e-mail was a scam.}
\end{quote}

European Portuguese; in Brazil people say \emph{o golpe}.

\begin{cousinweb}
\emph{burla} came from Spanish \emph{burla} \textquotedblleft{}mockery, a trick\textquotedblright{} (EtymDB), which goes back to Late Latin \emph{burra} \textquotedblleft{}nonsense, trickery\textquotedblright{}.
\end{cousinweb}

\begin{grammarlens}[title={What you've built}]
One more word for telling, explaining and describing at B1.
\end{grammarlens}

\subsection*{Your turn}
\begin{itemize}
\raggedright
  \item \emph{Say it:} \emph{a burla}
  \item \emph{Say it:} \emph{a burla}, once more
  \item \emph{Say it:} a sentence of your own with \emph{a burla}
  \item \emph{Recall:} say the Portuguese for the break-in, then the Portuguese for the fraud, then say \emph{a burla} again
\end{itemize}

\begin{tcolorbox}[breakable,colback=teal!4,colframe=teal!35!black,title={Before you move on}]
What does a burla mean? (\textquotedblleft{}The scam, the swindle\textquotedblright{}.) Say it once more.
\end{tcolorbox}
\section[burlar]{burlar — to swindle, to con}

\label{lesson:PT-C289-burlar}



\begin{quote}
Before the new one: say the Portuguese for the fraud, then the Portuguese for the scam.
\end{quote}

\subsection*{You'll want to know: burlar}
\textbf{burlar} — \textquotedblleft{}to swindle, to con\textquotedblright{}.

\begin{quote}
Um homem burlou várias pessoas idosas.
\end{quote}

\begin{quote}
\emph{A man swindled several elderly people.}
\end{quote}

In Brazil it means rather ‘to get round’ a rule: \emph{burlar a lei}.

\begin{cousinweb}
\emph{burlar} is built on \emph{burla} \textquotedblleft{}a trick, a scam\textquotedblright{} (EtymDB), from Spanish \emph{burla} \textquotedblleft{}mockery\textquotedblright{}.
\end{cousinweb}

\begin{grammarlens}[title={What you've built}]
One more word for telling, explaining and describing at B1.
\end{grammarlens}

\subsection*{Your turn}
\begin{itemize}
\raggedright
  \item \emph{Say it:} \emph{burlar}
  \item \emph{Say it:} \emph{burlar}, once more
  \item \emph{Say it:} a sentence of your own with \emph{burlar}
  \item \emph{Recall:} say the Portuguese for the fraud, then the Portuguese for the scam, then say \emph{burlar} again
\end{itemize}

\begin{tcolorbox}[breakable,colback=teal!4,colframe=teal!35!black,title={Before you move on}]
What does burlar mean? (\textquotedblleft{}To swindle, to con\textquotedblright{}.) Say it once more.
\end{tcolorbox}
\section[o burlão]{o burlão — the con man, the swindler}

\label{lesson:PT-C289-burlao}



\begin{quote}
Before the new one: say the Portuguese for the scam, then the Portuguese for to swindle.
\end{quote}

\subsection*{You'll want to know: o burlão}
\textbf{o burlão} — \textquotedblleft{}the con man, the swindler\textquotedblright{}.

\begin{quote}
O burlão prometia dinheiro fácil.
\end{quote}

\begin{quote}
\emph{The con man promised easy money.}
\end{quote}

European Portuguese; a woman is \emph{a burlona}.

\begin{cousinweb}
\emph{burlão} is built on \emph{burla} \textquotedblleft{}a scam\textquotedblright{} with the ending \emph{-ão} for a person; \emph{burla} came from Spanish (EtymDB).
\end{cousinweb}

\begin{grammarlens}[title={What you've built}]
One more word for telling, explaining and describing at B1.
\end{grammarlens}

\subsection*{Your turn}
\begin{itemize}
\raggedright
  \item \emph{Say it:} \emph{o burlão}
  \item \emph{Say it:} \emph{o burlão}, once more
  \item \emph{Say it:} a sentence of your own with \emph{o burlão}
  \item \emph{Recall:} say the Portuguese for the scam, then the Portuguese for to swindle, then say \emph{o burlão} again
\end{itemize}

\begin{tcolorbox}[breakable,colback=teal!4,colframe=teal!35!black,title={Before you move on}]
What does o burlão mean? (\textquotedblleft{}The con man, the swindler\textquotedblright{}.) Say it once more.
\end{tcolorbox}
\section[o cúmplice]{o cúmplice — the accomplice}

\label{lesson:PT-C289-cumplice}



\begin{quote}
Before the new one: say the Portuguese for to swindle, then the Portuguese for the con man.
\end{quote}

\subsection*{You'll want to know: o cúmplice}
\textbf{o cúmplice} — \textquotedblleft{}the accomplice\textquotedblright{}.

\begin{quote}
O burlão tinha um cúmplice dentro do banco.
\end{quote}

\begin{quote}
\emph{The con man had an accomplice inside the bank.}
\end{quote}

The same form for a woman: \emph{a cúmplice}.

\begin{cousinweb}
\emph{cúmplice} is a learned borrowing from Late Latin \emph{complex} \textquotedblleft{}closely joined, a partner\textquotedblright{} (EtymDB), from \emph{complicāre} \textquotedblleft{}to fold together\textquotedblright{}.
\end{cousinweb}

\begin{grammarlens}[title={What you've built}]
One more word for telling, explaining and describing at B1.
\end{grammarlens}

\subsection*{Your turn}
\begin{itemize}
\raggedright
  \item \emph{Say it:} \emph{o cúmplice}
  \item \emph{Say it:} \emph{o cúmplice}, once more
  \item \emph{Say it:} a sentence of your own with \emph{o cúmplice}
  \item \emph{Recall:} say the Portuguese for to swindle, then the Portuguese for the con man, then say \emph{o cúmplice} again
\end{itemize}

\begin{tcolorbox}[breakable,colback=teal!4,colframe=teal!35!black,title={Before you move on}]
What does o cúmplice mean? (\textquotedblleft{}The accomplice\textquotedblright{}.) Say it once more.
\end{tcolorbox}
